\begin{afterword}
\summaryhead

The post gives two cases of people who would not give up. Casey Serin, a 24-year-old web programmer with no experience in real estate, lied on mortgage applications to buy eight houses in several states, took cash out of the loans, and now owes \$2.2 million. He still will not declare bankruptcy or get a job, keeps paying for real-estate seminars, and tried to buy a ninth house.

The second case is Long-Term Capital Management, whose partners included the Nobel laureates Merton and Scholes. After three years of large profits, its edge vanished as others copied its trades. It borrowed more and more to exploit smaller and smaller margins, and by the time things went wrong it had \$4.72 billion of equity, \$124.5 billion of borrowing and \$1.25 trillion in derivatives.

The author concludes that each profession has its own skills, but that not turning little mistakes into big ones seems much the same art everywhere, and one of its keys is to be ready to admit you lost.

\responsehead

The post's facts hold up. Contemporary reports confirm Serin's age, job, eight houses in four states and \$2.2 million of debt, and the LTCM figures match those Wikipedia gives for the start of 1998. On LTCM the record supports the post's reading more than the post shows: in the crisis of 1998 the firm's founder ``still believed that their trades would pay off if they could simply weather the storm.''

The thesis is broader than its evidence. It is offered with ``it seems to me,'' which is fair, but it claims an art common to ``hedge funds or romance'' and to every profession. Both examples are debt-financed bets on markets, and the other professions get no example.

The main gap is in the rule itself. ``Be ready to admit you lost'' assumes that one can tell a loss from a setback, and the post gives no way to do so. Its second case shows why one is needed. On Wikipedia's account, the positions LTCM could not hold were taken over by a group of banks and ``eventually liquidated at a small profit to the rescuers.'' If that is right, the partners' belief in their trades was not simply mistaken. My inference is that what sank the fund was the borrowing, which left it unable to wait. The lesson the case supports is about leverage, which the post mentions, more than about hope.

The post introduces LTCM through its two Nobel laureates, then speaks of the firm. The laureates were partners. The firm was founded and run by John Meriwether, and a Federal Reserve history reports that Scholes was the partner most uneasy about the fund's move into new markets in 1997.

In short: two accurately reported cases of people who kept going, and a rule, admit you lost, that gives no way to tell a loss from a setback, drawn in part from a case whose trades, on one account, eventually made a small profit for the banks that took them over.
\end{afterword}
